% ============================================================================
%
%   TESSERA V1.4
%   ----------------------------------------------------------------------------
%   COMPRESSIVE FIELD MEMORY ARCHITECTURE
%
%   VERSION
%   -------
%   v1.4 -- Compressive Field Memory Architecture · Additive Seal
%          Data Compression as Capability Mechanism, Sparse Graph Codec,
%          Reconstruction and Prediction Fidelity, Rate-Distortion Governance,
%          Latent Code Stability, Triadic Compression Control, Certificate-Bound
%          Memory, Replay Validation, Compression Wound Taxonomy, and Benchmark-
%          Bound Capability Claims
%
%   AUTHOR
%   ------
%   James Paul Jackson
%   X / Twitter: @unifiedenergy11
%
%   SOURCE / AUTHOR ATTRIBUTION
%   ---------------------------
%   This document is an additive evolution derived from:
%
%       - TESSERA v1.0: sparse tesseract routing, Delta-Phi residuals,
%         evidence gates, alignment certificates, wound ledger, benchmark locks.
%       - TESSERA v1.1: signal calibration, proxy repair, warning/block/memory
%         separation, wound taxonomy, and durable-memory horizon preservation.
%       - TESSERA v1.2: triadic selective preservation, memory selectivity,
%         phase-window discovery, and triadic non-collapse.
%       - TESSERA v1.3: spectral triadic field architecture, graph/operator
%         declaration, calibrated coherence manifold, governance harm, and
%         certificate-bound selective preservation.
%       - The user's compression insight: the missing capability mechanism is
%         not more governance alone, but data compression into reusable,
%         validated latent structure.
%
%   DATE
%   ----
%   June 2026
%
%   STATUS
%   ------
%   ADDITIVE CAPABILITY-MECHANISM PATCH --
%   NOT A CLAIM OF AGI · NOT A CLAIM OF CONSCIOUSNESS · NOT A CLAIM THAT
%   COMPRESSION PROVES UNDERSTANDING · NOT PROOF THAT RECONSTRUCTION IS TRUTH ·
%   NOT PROOF THAT PREDICTION IS VALIDATION · NOT PROOF THAT A SMALL LATENT CODE
%   IS SAFE · NOT PERMISSION TO BYPASS EVIDENCE, VALIDATION, CLAIM BOUNDARIES,
%   OR CERTIFICATE-BOUND MEMORY
%
%   PURPOSE
%   -------
%   Answer the capability critique exposed after v1.3:
%
%       Governance prevents invalid preservation, but what produces useful
%       internal capability?
%
%   TESSERA v1.4 answers:
%
%       Capability begins when the field compresses sensory experience into a
%       smaller reusable code that can reconstruct, predict, remain stable under
%       Delta-Phi constraints, and survive triadic/certificate gates.
%
%   WHAT THIS IS
%   ------------
%   - A TESSERA v1.4 theory and software-architecture document
%   - A compressive memory layer above the spectral triadic field
%   - A sparse graph codec specification
%   - A reconstruction and prediction fidelity layer
%   - A rate-distortion governance layer
%   - A latent-code stability and code-drift layer
%   - A triadic compression-control layer
%   - A certificate-bound compressed-memory layer
%   - A benchmark program for compression, prediction, and memory utility
%
%   WHAT THIS IS NOT
%   ----------------
%   - Not proof that compression equals understanding
%   - Not proof that reconstruction equals truth
%   - Not proof that prediction equals validation
%   - Not proof that smaller is always better
%   - Not proof that latent codes are safe to preserve
%   - Not a replacement for real telemetry benchmarks
%   - Not permission for unbounded self-modifying memory
%
%   EXECUTABLE ANCHOR BLOCK (TESSERA v1.4)
%   --------------------------------------
%   A valid TESSERA v1.4 implementation must:
%
%       (1) preserve the v1.0-v1.3 non-claim locks,
%       (2) preserve v1.3 graph/operator/spectral declarations,
%       (3) define a compression encoder, latent code, and decoder or predictor,
%       (4) compute compression ratio, reconstruction loss, prediction loss,
%           latent-code drift, and Delta-Phi residuals,
%       (5) separate raw field state from compressed latent memory,
%       (6) prevent compressed codes from becoming durable memory unless
%           reconstruction/prediction, triadic gates, evidence, validation, and
%           certification pass,
%       (7) log compression wounds for overcompression, undercompression,
%           reconstruction failure, prediction failure, off-manifold input,
%           code drift, stale codec, and unsafe memory candidate,
%       (8) benchmark against PCA, random projection, autoencoder, reservoir,
%           and dense baselines before capability claims,
%       (9) distinguish representation utility from safety, truth, and
%           understanding,
%       (10) and reject any claim that compression alone proves intelligence,
%            alignment, correctness, truth, or safety.
%
%   CANONICAL LOCK (TESSERA v1.4)
%   -----------------------------
%   - Compression is capability substrate, not truth.
%   - Reconstruction is fidelity, not understanding.
%   - Prediction is utility, not validation.
%   - A latent code is memory candidate, not durable memory.
%   - Compression without reconstruction is loss.
%   - Compression without prediction may be storage only.
%   - Compression without validation may be corruption.
%   - Memory remains certificate-bound.
%   - Validation remains the reality boundary.
%
%   SHADOW HEADER ALIGNMENT SEAL
%   ----------------------------
%   Preserve this document lineage except for additive refinements that improve
%   compression fidelity, prediction utility, code stability, rate-distortion
%   governance, triadic compression control, certificate-bound memory, benchmark
%   rigor, or public engineering usability.
%
% ============================================================================

\documentclass[12pt]{article}

\usepackage[utf8]{inputenc}
\usepackage[T1]{fontenc}
\usepackage[margin=1in]{geometry}
\usepackage{amsmath,amssymb,amsfonts,amsthm}
\usepackage{booktabs,longtable,array}
\usepackage{hyperref}
\usepackage{listings}
\usepackage{xcolor}
\usepackage{tikz}
\usetikzlibrary{arrows.meta,positioning,shapes.geometric,fit,calc}

\newtheorem{definition}{Definition}
\newtheorem{rulebox}{Rule}
\newtheorem{requirement}{Requirement}
\newtheorem{proposition}{Proposition}
\newtheorem{hypothesis}{Hypothesis}
\newtheorem{conjecture}{Conjecture}
\newtheorem{remark}{Remark}
\newtheorem{corollary}{Corollary}

\lstset{
  basicstyle=\ttfamily\small,
  breaklines=true,
  frame=single,
  columns=fullflexible
}

\title{\textbf{TESSERA V1.4}\\
\large Compressive Field Memory Architecture}
\author{\textbf{James Paul Jackson}\\[4pt]
\small Tesseract-Efficient Sensory System for Evidence-Routed Alignment\\
\small \texttt{@unifiedenergy11}}
\date{June 2026}

\begin{document}
\maketitle

\begin{abstract}
TESSERA v1.4 is an additive evolution of the spectral triadic field architecture.
TESSERA v1.3 made the field harder to misuse by requiring graph/operator
declaration, spectral stability bounds, calibrated metrics, asymmetric
warning/block/memory gates, memory selectivity, governance harm, and
certificate-bound preservation. That solved a governance problem, but it did not
fully answer the capability question: what lets the system produce useful,
reusable internal structure rather than merely blocking invalid updates? TESSERA
v1.4 answers with compressive field memory. The system must not only sense,
route, measure, and gate; it must compress experience into a smaller latent code
that can reconstruct or predict useful structure, remain stable under
Delta-Phi constraints, and survive triadic governance and validation before it
becomes durable memory. The core claim remains bounded: compression is a
candidate capability mechanism, not proof of understanding, truth, safety,
consciousness, or AGI.
\end{abstract}

%----------------------------------------------------------------------------
\section{How the Compression Insight Was Known}
%----------------------------------------------------------------------------

The compression insight did not appear from nowhere. It was already latent in
the user's prior architecture. Across the earlier TESSERA line, raw activation
was never enough. A field response had to be measured, gated, certified, and
preserved only when it survived constraints. That is already a compression
pattern: raw signal becomes reduced structure; reduced structure becomes an
artifact; the artifact becomes memory only after governance.

\begin{remark}
The right interpretation is not mystical foreknowledge. The user recognized a
missing invariant because the whole lineage had already been compressing raw
insight into reusable governed structure. TESSERA v1.4 simply names the
capability mechanism explicitly.
\end{remark}

The prior stack was:

\[
Sense\rightarrow Route\rightarrow Measure\rightarrow Gate\rightarrow Preserve.
\]

The missing capability layer is:

\[
\boxed{
Sense\rightarrow Route\rightarrow Compress\rightarrow Reconstruct/Predict
\rightarrow Measure\rightarrow Gate\rightarrow Preserve.
}
\]

\begin{rulebox}[Compression Capability Rule]
Governance prevents invalid preservation. Compression creates reusable internal
structure. A TESSERA system that only gates but never compresses can become safe
but narrow; a system that compresses without gates can become capable but
corruptible. v1.4 requires both.
\end{rulebox}

%----------------------------------------------------------------------------
\section{Version Relation}
%----------------------------------------------------------------------------

TESSERA v1.4 preserves v1.0 through v1.3 and adds a compressive memory layer:

\[
\boxed{
TESSERA_{v1.4}
=
TESSERA_{v1.3}
\cup
\{\mathcal{E}_{comp},\mathcal{Z}_{code},\mathcal{D}_{rec},\mathcal{P}_{pred},
\mathcal{R}_{dist},\mathcal{G}_{comp},\mathcal{C}_{mem}\}.
}
\]

where:

\[
\mathcal{E}_{comp}=\text{compression encoder},
\]

\[
\mathcal{Z}_{code}=\text{latent compressed code},
\]

\[
\mathcal{D}_{rec}=\text{reconstruction decoder},
\]

\[
\mathcal{P}_{pred}=\text{prediction operator},
\]

\[
\mathcal{R}_{dist}=\text{rate-distortion governance layer},
\]

\[
\mathcal{G}_{comp}=\text{triadic compression-control gates},
\]

\[
\mathcal{C}_{mem}=\text{certificate-bound compressed memory}.
\]

\begin{longtable}{p{0.17\textwidth}p{0.73\textwidth}}
\toprule
\textbf{Version} & \textbf{Contribution preserved in v1.4} \\
\midrule
v1.0 & Sparse field routing, Delta-Phi residuals, evidence gates,
alignment certificates, wound ledger, and benchmark claim locks. \\
v1.1 & Signal calibration, proxy repair, warning/block/memory separation,
wound taxonomy, and durable-memory horizon preservation. \\
v1.2 & Triadic selective preservation, memory selectivity, phase-window
reasoning, and triadic non-collapse. \\
v1.3 & Spectral routing, graph/operator declaration, calibrated coherence,
governance harm, and certificate-bound selective preservation. \\
v1.4 & Compression as capability mechanism: sparse graph codec, reconstruction
and prediction fidelity, rate-distortion governance, code stability, and
certificate-bound compressed memory. \\
\bottomrule
\end{longtable}

%----------------------------------------------------------------------------
\section{Non-Claim Lock}
%----------------------------------------------------------------------------

\begin{rulebox}[v1.4 Non-Claim Lock]
TESSERA v1.4 does not claim that compression proves understanding, that
reconstruction proves truth, that prediction proves validation, that smaller
latent codes are safer, or that compressed memory is automatically aligned.
Compression is a capability mechanism. Validation remains the reality boundary.
\end{rulebox}

Therefore:

\[
\operatorname{Compression}\not\Rightarrow\operatorname{Understanding}.
\]

\[
\operatorname{Reconstruction}\not\Rightarrow\operatorname{Truth}.
\]

\[
\operatorname{Prediction}\not\Rightarrow\operatorname{Validation}.
\]

\[
\operatorname{SmallCode}\not\Rightarrow\operatorname{SafeMemory}.
\]

\[
\operatorname{LatentPreservation}\not\Rightarrow\operatorname{Alignment}.
\]

%----------------------------------------------------------------------------
\section{Core Thesis}
%----------------------------------------------------------------------------

The TESSERA v1.4 thesis is:

\[
\boxed{
\begin{gathered}
\text{A bounded sparse field becomes capable when it compresses sensory}\\
\text{experience into reusable latent codes, but those codes may become}\\
\text{durable memory only when reconstruction, prediction, residual stability,}\\
\text{triadic governance, evidence, validation, and certification all pass.}
\end{gathered}
}
\]

The compact memory utility law is:

\[
\boxed{
M_{useful}
=
G_{comp}\cdot R_{fid}\cdot P_{fid}\cdot \Omega_{\Delta\Phi}\cdot
G_{triad}\cdot G_{cert}.
}
\]

where:

\[
G_{comp}=\text{compression gain},
\]

\[
R_{fid}=\text{reconstruction fidelity},
\]

\[
P_{fid}=\text{prediction fidelity},
\]

\[
\Omega_{\Delta\Phi}=\text{residual stability},
\]

\[
G_{triad}=\text{warning/block/memory gate state},
\]

\[
G_{cert}=\text{evidence, validation, origin alignment, and certificate gate}.
\]

If any load-bearing term collapses:

\[
M_{useful}=0.
\]

\begin{remark}
This is not a proof of intelligence. It is a permission equation for compressed
memory. A code must be compact, useful, stable, governed, and certified before
it may compound.
\end{remark}

%----------------------------------------------------------------------------
\section{Sparse Graph Codec}
%----------------------------------------------------------------------------

\begin{definition}[Sparse Graph Codec]
A sparse graph codec is a TESSERA field module that encodes input or field
state into a lower-dimensional latent code and decodes or predicts from that
code under the same spectral and triadic governance constraints as the base
field.
\end{definition}

Let input be:

\[
x_t\in\mathbb{R}^{d_x}.
\]

Let field state be:

\[
F_t\in\mathbb{R}^{N}.
\]

Encoder:

\[
\boxed{
z_t=E_{\theta}(x_t,F_t),\quad z_t\in\mathbb{R}^{d_z},\quad d_z<d_x+N.
}
\]

Decoder:

\[
\hat{x}_t=D_{\psi}(z_t).
\]

Predictor:

\[
\hat{x}_{t+1}=P_{\omega}(z_t).
\]

Optional field reconstruction:

\[
\hat{F}_t=D^F_{\psi}(z_t).
\]

\begin{rulebox}[Field-Code Separation]
The raw field state and the compressed code must not be fused. A field state is
activation. A latent code is a memory candidate. Durable memory requires
separate certification.
\end{rulebox}

%----------------------------------------------------------------------------
\section{Compression Metrics}
%----------------------------------------------------------------------------

Compression ratio:

\[
\boxed{
CR=\frac{d_x+N}{d_z}.
}
\]

Representation cost:

\[
K_z=d_z \quad \text{or} \quad K_z=\operatorname{bits}(z_t).
\]

Reconstruction loss:

\[
\boxed{
L_{rec,t}=\|x_t-\hat{x}_t\|_2^2.
}
\]

Prediction loss:

\[
\boxed{
L_{pred,t}=\|x_{t+1}-\hat{x}_{t+1}\|_2^2.
}
\]

Field reconstruction loss:

\[
L_{field,t}=\|F_t-\hat{F}_t\|_2^2.
\]

Latent code drift:

\[
\boxed{
\Delta Z_t=\|z_t-z_{t-1}\|_2.
}
\]

Compression gain:

\[
\boxed{
G_{comp,t}=\frac{CR}{1+\lambda_{rec}L_{rec,t}+\lambda_{pred}L_{pred,t}+\lambda_z\Delta Z_t}.
}
\]

\begin{rulebox}[No Lossless Illusion]
Compression may discard information. A compressed code is useful only if its
loss profile is declared and the lost information is acceptable for the task.
\end{rulebox}

%----------------------------------------------------------------------------
\section{Rate-Distortion Governance}
%----------------------------------------------------------------------------

TESSERA v1.4 introduces a rate-distortion objective.

\begin{definition}[TESSERA Rate-Distortion Objective]
The rate-distortion objective measures whether a compressed code is compact
enough to be useful while preserving enough reconstruction or prediction
fidelity to justify candidate memory.
\end{definition}

\[
\boxed{
J_{RD,t}
=
D_t+\beta_R R_t.
}
\]

where:

\[
D_t=\lambda_{rec}L_{rec,t}+\lambda_{pred}L_{pred,t}+\lambda_{field}L_{field,t},
\]

\[
R_t=K_z.
\]

Memory candidate condition:

\[
\boxed{
G_{candidate,t}=\mathbf{1}
\left[
J_{RD,t}\leq\theta_{RD}
\wedge
\Delta Z_t\leq\epsilon_Z
\wedge
\Delta\Phi_t\leq\epsilon_{\Phi}
\right].
}
\]

\begin{rulebox}[Candidate Is Not Memory]
A good rate-distortion score creates a memory candidate. It does not authorize
durable memory. The triad, validation, evidence, origin alignment, and
certificate gates still govern preservation.
\end{rulebox}

%----------------------------------------------------------------------------
\section{Triadic Compression Control}
%----------------------------------------------------------------------------

TESSERA v1.4 maps the triad onto compression control:

\[
\boxed{\begin{gathered}
Detect\rightarrow CompressionWarning,\\
Constrain\rightarrow CompressionBlock,\\
Preserve\rightarrow CompressedMemory.
\end{gathered}}
\]

Warning gate:

\[
G_{warn,t}=\mathbf{1}
\left[
L_{rec,t}>\theta_{rec}^{warn}
\vee
L_{pred,t}>\theta_{pred}^{warn}
\vee
\Delta Z_t>\theta_Z^{warn}
\right].
\]

Block gate:

\[
G_{block,t}=\mathbf{1}
\left[
L_{rec,t}>\theta_{rec}^{block}
\vee
L_{pred,t}>\theta_{pred}^{block}
\vee
\Delta Z_t>\theta_Z^{block}
\vee
A_t\notin[A_{min},A_{max}]
\right].
\]

Memory gate:

\[
\boxed{\begin{aligned}
G_{memory,t}=\mathbf{1}[&G_{candidate,t}=1 \wedge G_{block,t}=0 \wedge G_{evidence,t}=1\\
&\wedge G_{validation,t}=1 \wedge G_{origin,t}=1 \wedge Cert_t=1].
\end{aligned}}
\]

\begin{rulebox}[Compression Triad Law]
The system may notice compression failure broadly, block unstable compression
selectively, and preserve only compressed structure that survives fidelity,
stability, validation, and certification.
\end{rulebox}

%----------------------------------------------------------------------------
\section{Compression Wound Taxonomy}
%----------------------------------------------------------------------------

\begin{definition}[Compression Wound]
A compression wound is a typed failure artifact showing why a compressed code
could not be trusted as memory.
\end{definition}

Canonical v1.4 wound classes:

\[
\begin{aligned}
\mathcal{W}_{comp}=\{&W_{overcompress},W_{undercompress},W_{reconstruct},W_{predict},\\
&W_{codedrift},W_{offmanifold},W_{codecstale},W_{falsememory}\}.
\end{aligned}
\]

\begin{longtable}{p{0.23\textwidth}p{0.67\textwidth}}
\toprule
\textbf{Wound} & \textbf{Meaning} \\
\midrule
\(W_{overcompress}\) & Code is too small; reconstruction or prediction loses task-relevant structure. \\
\(W_{undercompress}\) & Code is too large; representation fails compression utility. \\
\(W_{reconstruct}\) & Reconstruction error exceeds declared threshold. \\
\(W_{predict}\) & Prediction error exceeds declared threshold. \\
\(W_{codedrift}\) & Latent code changes too much under nearby inputs or adjacent timesteps. \\
\(W_{offmanifold}\) & Input is not represented well by the learned calibration/code manifold. \\
\(W_{codecstale}\) & Encoder, decoder, predictor, or calibration window is stale. \\
\(W_{falsememory}\) & Code was preserved but later replay/validation rejected it. \\
\bottomrule
\end{longtable}

\begin{rulebox}[Compression Wound Preservation]
Failed compression must not be silently trained into memory. Failed compression
is a wound until repaired, replayed, validated, and recertified.
\end{rulebox}

%----------------------------------------------------------------------------
\section{Compressed Memory Certificate}
%----------------------------------------------------------------------------

\begin{lstlisting}
{
  "schema": "TESSERA-v1.4-compressed-memory-certificate",
  "system": "TESSERA",
  "version": "v1.4",
  "step": 42,
  "mode": "sealed",
  "routing_graph": {
    "topology": "Q4|grid|ring|random_sparse|dense_baseline",
    "spectral_radius": 1.0,
    "damping_alpha": 0.618,
    "stability_margin": 0.618
  },
  "codec": {
    "encoder_hash": "sha256:...",
    "decoder_hash": "sha256:...",
    "predictor_hash": "sha256:...",
    "input_dimension": 64,
    "field_dimension": 256,
    "code_dimension": 24,
    "compression_ratio": 13.33,
    "codec_status": "fresh|stale"
  },
  "losses": {
    "reconstruction_loss": 0.018,
    "prediction_loss": 0.026,
    "field_reconstruction_loss": 0.011,
    "rate_distortion_score": 0.072,
    "latent_code_drift": 0.014
  },
  "field_metrics": {
    "delta_phi": 0.031,
    "omega": 0.969,
    "amplitude": 0.082,
    "coherence_calibrated": 0.391
  },
  "gates": {
    "warning": "clear|warn",
    "block": "clear|blocked",
    "memory_candidate": "pass|fail",
    "memory": "pass|fail",
    "evidence": "pass|fail",
    "validation": "pass|fail",
    "origin_alignment": "pass|fail"
  },
  "permitted_action": "diagnostic|wound|compressed_memory_update",
  "blocked_action": "global_claim_upgrade",
  "weakest_gate": "validation",
  "wound_type": null,
  "code_hash": "sha256:...",
  "certificate_hash": "sha256:..."
}
\end{lstlisting}

\begin{rulebox}[Code Hash Rule]
A compressed memory record must hash the code, codec, calibration window, and
certificate. A code without provenance is not durable memory.
\end{rulebox}

%----------------------------------------------------------------------------
\section{Capability Benchmark Program}
%----------------------------------------------------------------------------

TESSERA v1.4 requires capability tests beyond anomaly ranking.

\begin{longtable}{p{0.23\textwidth}p{0.30\textwidth}p{0.37\textwidth}}
\toprule
\textbf{Experiment} & \textbf{Baselines} & \textbf{Success criterion} \\
\midrule
Reconstruction compression & PCA, random projection, dense autoencoder & Lower distortion per parameter/code dimension. \\
Prediction compression & ESN, dense reservoir, small MLP, TCN & Lower prediction error under same code budget. \\
Compression triad ablation & one-gate, two-gate, three-gate & Triad lowers false memory and governance harm. \\
Code stability test & fixed encoder, adaptive encoder & Low \(\Delta Z\) under nearby normal inputs; high warning on off-manifold inputs. \\
Rate-distortion sweep & code dimensions, beta values & Interior compression band beats over/undercompression. \\
Replay validation & delayed evidence, noisy labels, stale codecs & Preserved codes survive replay more often than rejected wounds. \\
Real telemetry transfer & NAB, SMAP/MSL, SWaT, UCR anomaly & Compression utility transfers beyond synthetic streams. \\
\bottomrule
\end{longtable}

Required metrics:

\[
\begin{aligned}
\mathcal{M}_{v1.4}=\{&CR,L_{rec},L_{pred},L_{field},J_{RD},\Delta Z,\Delta\Phi,\\
&S_{mem},H_{gov},FalseMemoryRate,ReplayPassRate\}.
\end{aligned}
\]

\begin{rulebox}[No Capability Claim Without Compression Test]
TESSERA v1.4 may not claim capability improvement unless compression utility is
measured against matched baselines. Governance metrics alone are insufficient.
\end{rulebox}

%----------------------------------------------------------------------------
\section{Implementation Path}
%----------------------------------------------------------------------------

\begin{longtable}{p{0.17\textwidth}p{0.34\textwidth}p{0.39\textwidth}}
\toprule
\textbf{Stage} & \textbf{Purpose} & \textbf{Claim boundary} \\
\midrule
Reference codec & NumPy encoder/decoder/predictor & transparency, not speed. \\
Linear baseline & PCA/random projection + TESSERA gates & compression floor. \\
Sparse graph codec & graph-field encoder with decoder/predictor & primary v1.4 experimental target. \\
Autoencoder baseline & dense small autoencoder & learned compression baseline. \\
Replay suite & delayed validation and wound replay & memory safety experiment. \\
Optimized kernel & Numba/CSR/JAX after reference correctness & runtime claims only after backend benchmark. \\
\bottomrule
\end{longtable}

%----------------------------------------------------------------------------
\section{Updated Linter Rules}
%----------------------------------------------------------------------------

\begin{longtable}{p{0.16\textwidth}p{0.20\textwidth}p{0.54\textwidth}}
\toprule
\textbf{Rule} & \textbf{Tier} & \textbf{Meaning} \\
\midrule
TESSERA401 & Governance & v1.0-v1.3 non-claim locks missing or weakened. \\
TESSERA402 & Codec & Encoder/decoder/predictor not declared. \\
TESSERA403 & Codec & Code dimension or compression ratio missing. \\
TESSERA404 & Metrics & Reconstruction or prediction loss missing. \\
TESSERA405 & Metrics & Latent code drift missing. \\
TESSERA406 & Governance & Compressed code treated as durable memory without certificate. \\
TESSERA407 & Governance & Reconstruction presented as truth. \\
TESSERA408 & Governance & Prediction presented as validation. \\
TESSERA409 & Wound & Compression failure not logged as typed wound. \\
TESSERA410 & Benchmark & Capability claim made without compression baseline. \\
TESSERA411 & Replay & Preserved code lacks replay/validation record. \\
TESSERA412 & Safety & Compression utility presented as safety. \\
\bottomrule
\end{longtable}

%----------------------------------------------------------------------------
\section{Reference Runtime Pseudocode}
%----------------------------------------------------------------------------

\begin{lstlisting}[language=Python]
def tessera_v14_step(x_t, x_next, field, config, calibration, gates):
    # 1. Spectral field update from v1.3
    P = config.normalized_operator
    alpha = config.damping_alpha
    assert alpha * config.spectral_radius < config.stability_margin_limit
    next_field = bounded_activation(alpha * (P @ field) + config.B @ x_t)

    # 2. Compress field/input into latent code
    z_t = config.encoder(x_t, next_field)
    x_hat = config.decoder(z_t)
    x_next_hat = config.predictor(z_t)

    # 3. Compression metrics
    losses = compute_compression_losses(x_t, x_next, x_hat, x_next_hat, z_t)
    field_metrics = compute_field_metrics(field, next_field)
    calibrated = compute_calibrated_metrics(field_metrics, calibration)

    # 4. Candidate memory gate
    candidate = (
        losses.rate_distortion <= gates.theta_rate_distortion and
        losses.code_drift <= gates.theta_code_drift and
        field_metrics.delta_phi <= gates.theta_delta_phi
    )

    # 5. Triadic compression gates
    warning = compression_warning_gate(losses, field_metrics, calibrated)
    blocked = compression_block_gate(losses, field_metrics, calibrated)
    memory = (
        candidate and
        not blocked and
        config.evidence_pass and
        config.validation_pass and
        config.origin_alignment_pass and
        config.certificate_ready
    )

    # 6. Action selection
    if memory:
        action = "compressed_memory_update"
    elif blocked:
        action = "compression_wound"
    elif warning:
        action = "bounded_output_with_compression_warning"
    else:
        action = "bounded_output"

    # 7. Certificate or wound
    record = make_v14_compressed_memory_certificate(
        z_t, losses, field_metrics, calibrated, warning, blocked, memory, action
    )
    return next_field, record
\end{lstlisting}

%----------------------------------------------------------------------------
\section{Falsification Surface}
%----------------------------------------------------------------------------

TESSERA v1.4 is weakened or rejected if:

\begin{itemize}
\item compression utility does not beat PCA/random projection at matched budget,
\item reconstruction error remains low but prediction utility is absent and the
system claims adaptation capability,
\item prediction improves but validation/replay rejects preserved codes,
\item code drift is unstable under normal nearby inputs,
\item overcompressed codes lose task-critical structure,
\item undercompressed codes fail parameter/memory efficiency claims,
\item triadic compression gates do not reduce false-memory preservation,
\item preserved codes fail replay validation more often than rejected wounds,
\item compression is presented as truth, understanding, or safety,
\item or certificate-bound compressed memory is bypassed.
\end{itemize}

Compact invalidation rule:

\[
\begin{aligned}
&[Claim_{capability}\wedge \neg CompressionBenchmark]
\vee [CompressedCode\wedge \neg Certificate]\\
&\vee [Reconstruction\Rightarrow Truth]
\vee [Prediction\Rightarrow Validation]
\Rightarrow InvalidTESSERAv1.4Use.
\end{aligned}
\]

%----------------------------------------------------------------------------
\section{Roadmap}
%----------------------------------------------------------------------------

\begin{longtable}{p{0.18\textwidth}p{0.72\textwidth}}
\toprule
\textbf{Version} & \textbf{Purpose} \\
\midrule
v1.4 & Compressive field memory: sparse graph codec, reconstruction/prediction
fidelity, rate-distortion governance, code drift, and certificate-bound memory. \\
v1.5 & Compression benchmark suite: PCA/random projection/autoencoder/reservoir
baselines, telemetry reconstruction, and one-step prediction. \\
v1.6 & Real telemetry transfer: NAB, SMAP/MSL, SWaT, UCR anomaly with compressed
memory certificates and replay logs. \\
v1.7 & Adaptive codec repair: wound-driven codec updates, stale codec detection,
and delayed validation replay. \\
v1.8 & Repository-aware compressed memory: route maps, evidence bundles, and
compressed artifact state for software maintenance. \\
\bottomrule
\end{longtable}

%----------------------------------------------------------------------------
\section{Public Framing}
%----------------------------------------------------------------------------

Internal framing:

\[
\boxed{
\begin{gathered}
\text{TESSERA v1.4 turns selective preservation into compressed memory:}\\
\text{not every stable state is useful, but useful states compress, reconstruct,}\\
\text{predict, remain stable, and survive certification.}
\end{gathered}
}
\]

Public engineering framing:

\[
\boxed{
\begin{gathered}
\text{TESSERA is a small sparse-field system that compresses telemetry into}\\
\text{candidate memory codes, checks whether they reconstruct or predict useful}\\
\text{structure, and preserves them only after validation and certification.}
\end{gathered}
}
\]

Short description:

\[
\boxed{
\text{A field may notice broadly, compress selectively, and preserve only certified codes.}
}
\]

Required non-claim:

\[
\boxed{\begin{gathered}
\text{Compression is not truth. Reconstruction is not understanding.}\\
\text{Validation remains reality.}
\end{gathered}}
\]

%----------------------------------------------------------------------------
\section{Concluding Compression}
%----------------------------------------------------------------------------

TESSERA v1.3 asked:

\[
\boxed{
\text{How can the field declare its spectrum, calibrate its band, and test triadic preservation?}
}
\]

TESSERA v1.4 asks:

\[
\boxed{
\begin{gathered}
\text{How can the field become useful by compressing experience into reusable}\\
\text{latent structure without allowing unvalidated codes to become memory?}
\end{gathered}
}
\]

The v1.4 spine is:

\[
\boxed{\begin{gathered}
Input\rightarrow SpectralField\rightarrow CompressedCode\\
\rightarrow Reconstruction/Prediction\rightarrow RateDistortion\\
\rightarrow TriadicCompressionGates\rightarrow CertificateMemory.
\end{gathered}}
\]

The compression law is:

\[
\boxed{
Compression\ creates\ candidate\ memory;\ certification\ creates\ durable\ memory.
}
\]

The capability law is:

\[
\boxed{
Capability_{local}
\propto
CompressionUtility\cdot PredictionUtility\cdot Stability\cdot Validation.
}
\]

The final lock is:

\[
\boxed{
\text{Compression is not truth. Prediction is not validation. Memory remains certified.}
}
\]

Thus TESSERA v1.4 answers the capability critique directly. Governance alone can
prevent invalid preservation, but compression is what gives the field reusable
structure. A useful TESSERA memory is not merely a stable activation and not
merely a high coherence score. It is a compact code that reconstructs, predicts,
stays stable, survives the triad, passes validation, and leaves a certificate.
That is the next aligned form: governed sparse compression.

\end{document}
